\subsection{Cross-Design Test: Real Covariates in the ABM}\label{sec:robustness-crossdesign}

Section~\ref{sec:method-estimators} flags an unresolved confound. The two estimators differ at once in modeling paradigm (household choice versus loan-level survival) and in data source (synthetic population versus real Freddie Mac loan records), so a low ABM share against a high hazard-framework share is equally consistent with either explanation. To make partial progress on separating them, I re-ran the ABM's decision rule on structural covariates --- coupon, loan age, original LTV --- from the same 75,000-loan Freddie Mac sample as Section~\ref{sec:hazard}, replacing the log-normal, FRED-median-calibrated draws used elsewhere. Behavioral draws (mobility desire, loss aversion, patience, transaction cost) stay synthetic; the Freddie Mac data contain no analogue for them. The test therefore varies real structural heterogeneity (and, until the post-stratification below, book composition jointly), not a fully real population.

The ABM's mobility-scale parameter is calibrated to hold a real-world baseline turnover floor fixed (Appendix~\ref{sec:method-abm}). Swapping in a population with different LTV and coupon moments therefore raises a calibration question a naive covariate swap does not resolve on its own: re-derive the turnover-floor anchor on the new moments, or hold it fixed? I report both variants. The recalibrated variant applies Appendix~\ref{sec:method-frictions}'s recalibration rule, re-anchoring the mobility scale to reproduce the same 4--5\% turnover floor under the real population's moments. The frozen variant suspends that rule and leaves the original synthetic-population calibration unchanged.

Before running either variant I fixed the interpretive thresholds ex ante. A recovery above 50\% of the empirical benchmark would undercut the paradigm interpretation advanced in Sections~\ref{sec:abm}--\ref{sec:conclusion}. A recovery below 35\% would corroborate it, consistent with the synthetic-population ABM's own specifications. A recovery between 35\% and 50\% would be ambiguous.

The recalibrated variant recovers \$453.5 billion, or 59.3\% of the benchmark, above my ex-ante 50\% threshold. The frozen variant recovers \$159.5 billion, or 20.9\%, within the corroborating band. Both are single draws. A pre-committed fifty-seed sweep over both randomization layers (behavioral draws and covariate subsample; run \texttt{cross\_design\_seeds}) shows the threshold crossing is the estimator's, not the seed's. The recalibrated variant averages 60.2\%, standard deviation 3.3 points, 95\% band 55.1--66.0\%. The committed draw sits at the 48th percentile of its own distribution, and all fifty seeds classify as undercutting. The frozen variant averages 22.0\%. The band is wide enough to report --- 10.9 points against the 7.7-point disclosure threshold fixed before the run --- but too narrow to reach the calibration swing. The operative uncertainty here is therefore calibration choice, not seed noise, and it is now measured. The same sweep records each variant's deep-gap floor, and the two separate cleanly on it. The recalibrated CPR at an 8\% market rate averages 4.59\%, inside the observed 4--5\% turnover band on all fifty seeds; the frozen averages 7.05\%, outside it on all fifty (range 6.31--7.63\%). Applied consistently, the admissibility test I use on the externally anchored variant earlier in this subsection counts against the frozen leg, not the recalibrated one: the leg reported as within the corroborating band is the one that misses the turnover moment. I therefore read the 20.9\% as the less admissible of the two calibrations, not as corroboration, which strengthens the undercutting reading. One caveat on ``frozen''. Re-deriving that variant's mobility-scale anchor per seed gives mean 43{,}493, standard deviation 1{,}181, minimum 39{,}984 against the committed 43{,}883, the very top of its own range. The anchor is frozen by convention, not because the data pin it there. Table~\ref{tab:crossdesign} summarizes both against the synthetic-population control; Figure~\ref{fig:crossdesign} plots the paired comparison.

\begin{figure}[H]
\centering
\includegraphics[width=\textwidth]{fig5_cross_design}
\caption{The cross-design test as paired panels. Panel (a): dollar recovery for the two real-covariate variants against the pre-committed thresholds (above 50\% undercuts the paradigm interpretation; below 35\% corroborates it) and the synthetic-population control band; the recalibrated variant's 59.3\% crosses the undercutting threshold while the frozen variant's 20.9\% stays within the corroborating band. Panel (b): lag-0 path correlations, all negative, so neither variant reproduces the empirical path's timing regardless of calibration choice, and after Section~\ref{sec:pathb}'s detrending result no estimator in this paper does.}
\label{fig:crossdesign}
\end{figure}

\begin{table}[H]
\centering
\small
\begin{threeparttable}
\caption{Cross-design test results: real Freddie Mac structural covariates substituted into the ABM under two calibration choices, against the synthetic-population production ABM.}
\label{tab:crossdesign}
\begin{tabular}{@{}p{4.0cm}lll@{}}
\toprule
 & Recalibrated (primary) & Frozen (robustness) & Synthetic control \\
\midrule
Mobility-scale parameter & 36,086 & 43,883 & 43,883 \\
Trapped liquidity & \$453.5B (59.3\%; fifty-seed mean 60.2\%) & \$159.5B (20.9\%; 22.0\%) & \$91.0B (11.9\%) \\
Mean CPR & 8.14\% & 11.63\% & 11.68\% \\
CPR path corr. (lag 0) & $-0.336$ & $-0.311$ & $-0.318$ \\
Peak lag & 0 (no positive alignment) & 0 (no positive alignment) & 0 \\
\bottomrule
\end{tabular}
\begin{tablenotes}[flushleft]\footnotesize
\item Peak lag follows the max-$|r|$ convention of Table~\ref{tab:estimators}. Under the alternative most-positive-$r$ convention, the synthetic control's best alignment is $+0.195$ at lag $-3$, the ``$+0.19$ at the three-month lead'' of Section~\ref{sec:abm-results}; the two conventions summarize the same cross-correlation profile, and the substantive conclusion (no positive contemporaneous alignment, which, after Section~\ref{sec:pathb}'s detrending result, no estimator including Path B shows) holds under either.
\end{tablenotes}
\end{threeparttable}
\end{table}

By my own ex-ante criterion the primary recalibrated variant's 59.3\% recovery, 60.2\% averaged over fifty seeds, moves against the paradigm interpretation. With real structural covariates and a like-for-like turnover anchor, the household-choice ABM explains far more of the benchmark than the synthetic-population specifications of Sections~\ref{sec:abm} and~\ref{sec:robustness-scale} suggested. But the \$294 billion gap between the recalibrated and frozen variants is produced solely by how the turnover-floor anchor is set, covariates fixed. It indicates that the ABM's aggregate share is highly sensitive to a calibration choice orthogonal to data source. The test therefore cannot, on its own, cleanly attribute Sections~\ref{sec:abm} and~\ref{sec:hazard}'s contrast to modeling paradigm rather than to data source: recalibration choice does comparable work to whichever the test was designed to isolate.

I therefore withdraw the paradigm reading as a claim of this paper. It was pre-committed to a threshold before the run; the threshold was crossed on the committed draw, on all fifty seeds, and by a wider margin once the sample carries the book's own composition. A reading that survives only by narrowing, after the fact, what the threshold was taken to test is not one I will defend, and narrowing it is what the pre-commitment architecture exists to prevent. Two results the ABM sections do establish stand without it, and neither is a claim about modeling paradigm: no synthetic-population household-choice specification tested here explains the Federal Reserve's cap shortfall, and no estimator in this paper reproduces the monthly path once the window's shared trend is removed.

A pre-committed post-stratification separates composition from heterogeneity. I reweight the real-covariate population to the SOMA coupon-by-vintage-group composition, the book the benchmark is measured on. On that population I re-derive the mobility scale under the same floor rule. That moves the recalibrated variant from 59.3\% to 76.3\% and the frozen variant from 20.9\% to 12.6\%. The reweight uses importance weights, effective sample size 1{,}332 of 10{,}000, disclosed. A seed-42 re-draw within SOMA-weighted cells gives 70.6\% for the recalibrated variant; the frozen variant's re-draw gives 36.0\%. Composition mismatch therefore understated rather than manufactured the recalibrated recovery. At book composition, real-covariate household choice with a like-for-like floor anchor explains roughly seventy to seventy-six percent of the benchmark. The ex-ante undercutting verdict survives composition adjustment, and the calibration-choice spread widens to roughly \$488 billion. The reweight is a level exercise; the path-level failure of Table~\ref{tab:crossdesign} stands.

A third variant drops the turnover-floor anchor entirely and takes both author-calibrated behavioral elements from Path B's external elasticity source. The mobility-desire scale comes from \citepos{liebersohn2024} zero-gap moving level: a 1.5\% quarterly ZIP-code moving hazard at an essentially zero average rate gap, 5.9\% annualized. The wait-and-see freeze share comes from what their 6.5\%-per-quarter-per-100 bp mobility decline implies at the production 150 bp trigger: 18.3\% against the author-calibrated 20\%. The DTI wall and loss-aversion coefficient are already externally sourced. Recovery does rise without the floor recalibration, to 150.6\% of benchmark. But it rises only by predicting essentially zero prepayment at QT-window gaps. CPR at an 8\% market rate is 0.0\%, against observed 4--5\% discount-cohort turnover and the roughly 1.44--1.45\% quarterly moving hazard \citet{liebersohn2024} themselves observe at 2022--2023 gaps. The externally anchored gate structure can match the zero-gap moving level or the observed deep-gap turnover, not both. The exponential mobility-desire distribution that fits the former (scale 4{,}472 against the recalibrated 36{,}086) shuts moving down once the par-payoff penalty binds. I read that as clarifying the floor anchor's role: it keeps the ABM on observed turnover data rather than suppressing its rate response. The anchored dimension is never claimed as fit, since matching the floor is tautological by construction (Appendix~\ref{sec:method-frictions}). What the variant measures is what that tautology purchases: empirical admissibility at deep gaps, another instance of this subsection's deflate-as-well-as-inflate indeterminacy (same harness). Scored against the pre-committed bands, as its own specification required, its 150.6\% classifies as undercutting the paradigm reading, and the frozen artifact records that classification. I do not read it as informative about the paradigm question, and I label that reading post hoc. The pre-run specification committed only to reporting the deep-gap floor diagnostic ``alongside whatever the recovery number is, whichever direction it moves''; it did not commit to removing the recovery from the band when the diagnostic fails. The admissibility precondition---an empirically admissible deep-gap turnover level before a recovery is interpretable---was formulated after the result was seen. I report it here as a labelled post-run qualification, not a pre-registered exclusion. A reader who declines it should score the variant at 150.6\% against the 50\% band. That strengthens the undercutting reading this subsection concedes, and no headline quantity turns on the choice.

I ran a minimal simulated-minimum-distance calibration to reconcile the two external moments: a joint two-parameter fit of the mobility-desire scale plus a forced-move point mass added to the structural family. Moments, admissibility bands, distance, and verdict vocabulary were fixed before execution, and a parity gate replayed the committed 59.303\% first. The verdict is joint infeasibility, now with two degrees of freedom, not one. At the distance-minimizing point (scale 1{,}000, mass 3.9\%) the zero-gap moving level overshoots its band, 0.0644 against $[0.055, 0.063]$, just as the deep-gap turnover reaches its own, 0.0394 against $[0.040, 0.050]$. The forced mass lifts both together at the scale's floor. The near-fit recovers 107.1\% of benchmark through the same harness, far above the pre-committed cross-design bands: external anchoring pushes this family past the benchmark, not toward it. I deliberately do not fit QT-window CPR moments. The recovery statistic is a functional of that path, so fitting it would make the test a fitted quantity---the overfitting the pre-commitment architecture exists to prevent. Hazard-side headline quantities are not functions of ABM calibration and cannot move under any outcome of this run.

Both variants retain the ABM's central path-level failure: the monthly CPR path correlation stays negative at lag 0 ($-0.336$ and $-0.311$). Every estimator's does, once the window's shared trend is removed (Section~\ref{sec:pathb}). Path B's distinction is level accuracy, not timing. The timing critique of Section~\ref{sec:abm} therefore survives whichever calibration variant is preferred. The dollar-recovery critique does not straightforwardly survive under the recalibrated variant.
